Flow categories originally emerged as a categorical bookkeeping mechanism to organize the geometric and analytic data extracted from Hamiltonian systems on a manifold \cite{CJS95}. In recent years, flow categories and their stable homotopy types have become central tools for lifting invariants in Floer, contact, and Khovanov homology into the realm of stable homotopy theory. A natural question arises: how can we systematically compute the generalized homology of these realization spectra directly from the underlying geometric flow data?   

Here is an intuitive description of the main theorem of the paper using a toy model coming from Morse theory.
Let $\Man$ be a closed manifold and $f: \Man \to \bbR$ a Smale Morse function. Denote by $\mathrm{Crit}(f)$ the set of critical points of $f$, and by $\ind: \mathrm{Crit}(f) \to \bbZ$ the index function. For simplicity, assume that there is one critical point for each index $i \in \{0, \dots, n\}$, and denote the critical points by $x_i$ where $\ind(x_i) = i$.
The first result is:

There is a spectral sequence 
\begin{equation*}
    \pg{E}{1}{n}{p} = \Omega^{\fr}_{n-p} \implies \pi_n \Sinf \Man.
\end{equation*}
where $\Omega^{\fr}_*$ is the framed bordism ring, and $\Sinf \Man$ is the stable homotopy type of  $\Man$.

Let $x_j,x_i$ be two critical points of $f$ with $j > i$. We denote by $\flowM(x_j,x_i)$ the compactification of the moduli space of flow lines from $x_j$ to $x_i$, up to reparametrization. Recall that $\flowM(x_j,x_i)$ is a framed manifold with corners of dimension $\ind(x_j) - \ind(x_i) - 1$, and that the boundary structure is given by
\begin{equation*}
    \partial \flowM(x_j,x_i) = \bigcup_{j > r > i} \flowM(x_j,x_r) \times \flowM(x_r,x_i).
\end{equation*}

Let $W_p$ be a representative of a bordism class in $\pg{E}{1}{n}{p} = \Omega^{\fr}_{n-p}$.
We have that $W_p$ survives to the $(r+1)$-th page (also denoted by $W_p \in Z_{r+1}$) if and only if there is a family of framed manifolds $\{W_s\}$ for $s \in \{p-1, \dots, p-r\}$ of dimension $n-s$ such that the corner structure is given by
\begin{equation*}
    \partial W_s = \bigcup_{p \geq t > s} W_t \times \flowM(x_t,x_s).
\end{equation*}

We can then glue the collection $W_s \times \flowM(x_s,x_{p-(r+1)})$ together.
Let $t,u \in \bbZ$ be such that $p \geq t > u \geq p-(r+1)$. By the definition of the boundary structure of $W_u$ and $\flowM(x_t,x_{p-(r+1)})$, we have that 
\begin{equation*}
    W_t \times \flowM(x_t,x_{p-(r+1)}) \hookleftarrow W_t \times \flowM(x_t,x_u) \times \flowM(x_u,x_{p-(r+1)}) \hookrightarrow W_u \times \flowM(x_u,x_{p-(r+1)})
\end{equation*}
and these boundary components are shared only by these two manifolds of the collection $W_s \times \flowM(x_s,x_{p-(r+1)})$.
On the other hand, all the boundary components of $W_s \times \flowM(x_s,x_{p-(r+1)})$ are of this form, so we can glue the collection $W_s \times \flowM(x_s,x_{p-(r+1)})$ together to obtain a framed closed manifold which we denote by
\begin{equation*}
    \flowM(-,x_{p-(r+1)}) \circ W \coloneqq \bigcup_{s \in \{p,\dots, p-r\}} \flowM(x_s,x_{p-(r+1)}) \times W_s,
\end{equation*}
which is an analog of the Toda bracket construction. Note that 
\begin{equation*}
    \dim(\flowM(x_s,x_{p-(r+1)}) \circ W) = \dim(W_s) + \dim(\flowM(x_s,x_{p-(r+1)})) = n-s + \big(s - (p-(r+1)) - 1 \big) = n-p+r.
\end{equation*}
The last part of the theorem states that if $W_p \in \pg{E}{1}{n}{p} = \Omega^{\fr}_{n-p}$ survives to $Z_{r+1}$, then the $r+!$ differential on $Z_{r+1}$ is given by the bordism class of the Toda manifold: 
\begin{equation*}
    d_{r+1}(W_p) = [\flowM(-,x_{p-(r+1)})\circ W] \in \Omega^{\fr}_{n-p+r} = \pg{E}{1}{n-1}{p-(r+1)}.
\end{equation*}
See \cref{fig:ffCat} for an example of the structure of the family $W_s$ and the flow manifolds,
\cref{fig:k-boundary} for the gluing diagram of the Toda manifold, and \cref{fig:differential} for the resulting closed manifold.
